% Vietnamese translation for OI Public Library
% Source-File: IOI中国国家候选队论文/2006/陈启峰/附件/消防站解题报告.doc
\documentclass[11pt]{article}
\usepackage{oipl-vn}

\title{Lời giải bài ``Trạm cứu hỏa''}
\author{Chen Qifeng}
\date{}

\begin{document}
\maketitle

\origtitle{Fire Station solution report}
\sourcepath{IOI 2006 / Chen Qifeng / attachments / Fire Station solution report.doc}

\section{Mô tả bài toán}

Quốc gia \(Z\) có \(n\) thành phố, đánh số từ \(1\) đến \(n\). Giữa các
thành phố có đường cao tốc, và giữa hai thành phố bất kỳ có đúng một đường đi.
Các đường cao tốc có thể có độ dài khác nhau.

Gần đây quốc gia \(Z\) thường xảy ra hỏa hoạn, nên chính quyền quyết định xây
một số trạm cứu hỏa. Nếu xây trạm cứu hỏa tại thành phố \(k\), chi phí phải
trả là \(W_k\). Nếu thành phố \(k\) không có trạm cứu hỏa, khoảng cách từ nó
đến trạm cứu hỏa gần nhất không được vượt quá \(D_k\). Mỗi thành phố, trong
điều kiện khoảng cách không vượt quá \(D_k\), phải chọn trạm cứu hỏa gần nhất
làm trạm phụ trách.

Để tiết kiệm, cần xây một số trạm cứu hỏa với tổng chi phí nhỏ nhất sao cho
mọi yêu cầu trên đều được thỏa mãn.

\section{Mô hình toán học}

Xem các thành phố là đỉnh, đường cao tốc là cạnh, độ dài đường là trọng số.
Tính chất ``giữa hai thành phố bất kỳ có đúng một đường đi'' cho thấy đồ thị
này là một cây.

Gọi \(\operatorname{dist}(u,v)\) là khoảng cách giữa hai đỉnh \(u\) và \(v\).
Ta cần tìm một dãy nhị phân \(x_1,x_2,\ldots,x_n\), trong đó \(x_i=1\) nghĩa
là xây trạm tại đỉnh \(i\), sao cho với mọi \(u\) đều tồn tại một đỉnh \(v\)
thỏa
\[
  x_v=1,\qquad \operatorname{dist}(u,v)\le D_u,
\]
và hàm mục tiêu
\[
  \sum_{i=1}^{n} W_i x_i
\]
đạt giá trị nhỏ nhất.

\section{Phân tích mô hình thuật toán}

Vì bài toán liên quan đến khoảng cách trên đồ thị, phản ứng đầu tiên là thử
các thuật toán đồ thị. Tuy nhiên sau nhiều cách thử, hướng đó không đem lại
một công thức hiệu quả.

Điểm đặc biệt cần khai thác là đồ thị là cây, và mục tiêu là cực tiểu hóa chi
phí. Từ hai đặc điểm này có thể gần như xác định rằng lời giải nên là quy
hoạch động trên cây.

\section{Mâu thuẫn khi xác định trạng thái}

Trong quy hoạch động, đặc biệt là quy hoạch động trên cây, việc xác định trạng
thái là phần quan trọng nhất. Thông thường trạng thái sẽ có một tham số \(i\),
biểu diễn đối tượng đang xét là cây con gốc \(i\).

Nếu chỉ dùng một trạng thái \(f_i\) để biểu diễn chi phí nhỏ nhất khi xây các
trạm trong cây con gốc \(i\), sao cho mọi đỉnh trong cây con đều cách trạm gần
nhất không quá giới hạn \(D\), thì không thể viết được chuyển tiếp. Lý do là
trạng thái này không phản ánh trạm nào đã được xây, cũng không cho biết chi
tiết về trạm gần nhất của \(i\).

Vì vậy ta muốn thêm một tham số nữa. Tham số này có thể là khoảng cách từ
\(i\) đến trạm gần nhất, cũng có thể là số hiệu trạm gần nhất của \(i\), hoặc
số hiệu trạm gần nhất nằm trong cây con, hoặc số hiệu trạm gần nhất nằm ngoài
cây con. Nhưng nếu giữ nguyên yêu cầu ``phải chọn trạm gần nhất'', các cách
định nghĩa đó đều không dẫn đến chuyển tiếp rõ ràng. Nguyên nhân chính là khi
ghép các cây con, ta không bảo đảm được thông tin ``gần nhất'' trong trạng
thái vẫn nhất quán.

\section{Dùng phương pháp nới điều kiện}

Rào cản lớn nhất là điều kiện nghiêm ngặt: mỗi thành phố, trong phạm vi cho
phép, bắt buộc phải chọn trạm cứu hỏa gần nhất làm trạm phụ trách. Thực ra để
bảo đảm tính hợp lệ, ta không cần biết trạm gần nhất là trạm nào; chỉ cần bảo
đảm trong khoảng cách \(D_i\) có ít nhất một trạm.

Vì thế ta nới điều kiện thành:

\begin{quote}
Mỗi thành phố, trong phạm vi khoảng cách không vượt quá \(D_i\), có thể chọn
bất kỳ một trạm cứu hỏa nào làm trạm phụ trách.
\end{quote}

Giá trị tối ưu sau khi nới điều kiện không đổi. Thật vậy, trong mọi phương án
xây trạm hợp lệ theo điều kiện nới lỏng, nếu một thành phố đang chọn một trạm
không gần nhất, ta đổi nó sang trạm gần nhất thì khoảng cách chỉ giảm, chi phí
xây trạm không đổi, và phương án vẫn hợp lệ. Do đó luôn tồn tại một phương án
tối ưu thỏa điều kiện ban đầu.

\section{Trạng thái sau khi nới điều kiện}

Sau khi bỏ chữ ``gần nhất'' khỏi định nghĩa, ta có thể dùng trạng thái hai
tham số. Gốc cây tại một đỉnh tùy ý. Ký hiệu \(T_i\) là cây con gốc \(i\).

Định nghĩa \(F(i,j)\) là chi phí nhỏ nhất thỏa các điều kiện sau:

\begin{itemize}
  \item xây một số trạm trong cây con \(T_i\);
  \item tại đỉnh \(j\) phải có một trạm cứu hỏa;
  \item mọi đỉnh trong \(T_i\), trong giới hạn khoảng cách của nó, chọn một
  trạm nằm trong \(T_i\) hoặc trạm tại \(j\) làm trạm phụ trách;
  \item riêng đỉnh \(i\) bắt buộc chọn trạm tại \(j\) làm trạm phụ trách;
  \item nếu \(j\notin T_i\), chi phí xây trạm tại \(j\) không được tính trong
  \(F(i,j)\).
\end{itemize}

Nhờ định nghĩa này, trạng thái chỉ ghi lại trạm phụ trách của gốc cây con,
không cần biết nó có phải trạm gần nhất hay không.

\section{Thêm điều kiện để có chuyển tiếp}

Sau khi nới điều kiện, trạng thái đã rõ hơn, nhưng nếu để mỗi thành phố chọn
trạm phụ trách tùy ý thì vẫn còn thiếu cấu trúc để chuyển tiếp. Quy hoạch động
cần thứ tự topo và tính không hậu nghiệm, nên ta thêm một điều kiện tự nhiên:

\begin{quote}
Nếu thành phố \(u\) chọn trạm tại thành phố \(v\), và đường đi từ \(u\) đến
\(v\) là \(u=p_0,p_1,\ldots,p_t=v\), thì mọi đỉnh \(p_r\) trên đường đi cũng
chọn trạm tại \(v\) làm trạm phụ trách.
\end{quote}

Cần chứng minh luôn có một phương án tối ưu thỏa điều kiện này.

\paragraph{Chứng minh.}
Lấy một phương án tối ưu bất kỳ, tức là một dãy nhị phân xác định các đỉnh có
trạm cứu hỏa. Thêm một đỉnh ảo \(s\), nối \(s\) với mọi đỉnh có trạm bằng cạnh
trọng số \(0\). Chạy Dijkstra từ \(s\) và ghi lại đỉnh trước trên đường đi
ngắn nhất.

Với mỗi đỉnh có trạm, chọn chính trạm của nó làm trạm phụ trách. Với mỗi đỉnh
không có trạm, chọn cùng trạm phụ trách với đỉnh trước của nó trên cây đường
đi ngắn nhất. Khi đó trên đường từ một đỉnh về trạm phụ trách, mọi đỉnh trung
gian đều chọn cùng trạm đó. Đồng thời vì mỗi đỉnh được gán đến một trạm gần
nhất trong phương án đã xây, khoảng cách đến trạm phụ trách không vượt quá
\(D_i\). Tập trạm được xây không thay đổi, nên tổng chi phí vẫn tối ưu.
Chứng minh hoàn tất.

Như vậy điều kiện bổ sung là hợp lệ và có thể dùng trong chuyển tiếp.

\section{Chuyển tiếp quy hoạch động}

Đặt
\[
  G(i)=\min_{j} F(i,j),
\]
nghĩa là chi phí nhỏ nhất để xây trạm hợp lệ trong cây con \(T_i\), khi mọi
đỉnh trong cây con đều chọn một trạm nằm trong cây con làm trạm phụ trách.

Trước hết, nếu
\[
  \operatorname{dist}(i,j)>D_i,
\]
thì \(F(i,j)=+\infty\), vì đỉnh \(i\) không thể chọn trạm \(j\).

Ngược lại, xét ba trường hợp.

\paragraph{Trạm \(j\) nằm ngoài \(T_i\).}
Với mỗi con \(v\) của \(i\), cây con \(T_v\) có hai lựa chọn: tự giải hoàn
toàn bên trong cây con với chi phí \(G(v)\), hoặc để gốc \(v\) chọn trạm ngoài
cây con là \(j\), với chi phí \(F(v,j)\). Do đó
\[
  F(i,j)=\sum_{v\in\operatorname{child}(i)}
  \min\{G(v),F(v,j)\}.
\]

\paragraph{\(j=i\).}
Ta phải trả thêm chi phí xây trạm tại \(i\). Các cây con vẫn có hai lựa chọn
như trên, nhưng trạm ngoài của chúng chính là \(i\):
\[
  F(i,i)=W_i+\sum_{v\in\operatorname{child}(i)}
  \min\{G(v),F(v,i)\}.
\]

\paragraph{\(j\in T_i\) và \(j\ne i\).}
Khi đó \(j\) nằm trong đúng một cây con \(T_c\) của một con \(c\) của \(i\).
Cây con \(T_c\) phải dùng trạng thái \(F(c,j)\). Với mọi cây con còn lại,
lựa chọn giống trường hợp trạm nằm ngoài cây con:
\[
  F(i,j)=F(c,j)+
  \sum_{\substack{v\in\operatorname{child}(i)\\v\ne c}}
  \min\{G(v),F(v,j)\}.
\]

Sau khi tính các \(F(i,j)\), lấy \(G(i)=\min_j F(i,j)\). Đáp án cuối cùng là
\(G(r)\), với \(r\) là gốc đã chọn của cây.

Số trạng thái là \(O(n^2)\). Nếu cài đặt trực tiếp các tổng theo công thức
trên, thời gian có thể là \(O(n^3)\) và bộ nhớ là \(O(n^2)\). Với việc tiền xử
lý khoảng cách trên cây và tối ưu các tổng theo cây con, có thể giảm hằng số
đáng kể, nhưng ý chính của lời giải nằm ở cách định nghĩa trạng thái.

\section{Kết luận}

Hai phương pháp ``nới điều kiện'' và ``thêm ràng buộc'' không mâu thuẫn với
nhau; ngược lại, chúng thường bổ sung cho nhau. Khi điều kiện quá nghiêm ngặt
khiến không thể định nghĩa trạng thái, ta nới điều kiện để nhìn ra cấu trúc.
Khi điều kiện sau đó quá rộng và thiếu hướng chuyển tiếp, ta thêm một ràng
buộc có thể chứng minh là vẫn giữ được nghiệm tối ưu.

Trong bài này, nới điều kiện giúp xác định trạng thái, còn thêm ràng buộc giúp
xác định công thức chuyển tiếp. Đó chính là cách ``một căng một chùng'' trong
thiết kế thuật toán.

\end{document}
